%% Положения, выносимые на защиту. ЕДИНСТВЕННЫЙ источник текста положений:
%% подключается общей характеристикой (common/characteristic.tex -> диссертация и автореферат)
%% и презентацией (Presentation/styles.tex). Правки вносить только здесь.
%%
%% ВНИМАНИЕ: файл используется и в тексте, и на слайдах, поэтому ссылки \cref/\ref
%% здесь недопустимы; указания на главы и разделы приводятся текстом.

% Историческая справка: положение о монокулярном алгоритме L-Shape исключено из
% выносимых на защиту (коммит 460959f, 14.09.2026); его результаты сохранены в
% главе 2 как промежуточные. Прежний текст:
% Алгоритм вписывания проекции объекта на виде сверху в ориентированную ограничивающую
% рамку, определяющий её ориентацию методом наименьших квадратов по модели L-образной
% формы видимой границы ТС, превосходит известные алгоритмы вписывания, не использующие
% обучаемых моделей, на тех же входных данных: по среднему коэффициенту Жаккара
% ($0,337$ против $0,328$ у лучшего из них), по всем квартилям распределения и при
% любом ограничении рабочей дальности.

\newcommand{\defposA}{Алгоритм оценки положения окружающих ТС на основе комплексирования визуальных и лидарных данных, выделяющий точки поверхности движения оценкой её формы регрессией на гауссовских процессах вместо предположения о плоскости, даёт наибольший средний коэффициент Жаккара среди алгоритмов, не использующих обучаемых преобразований в вид сверху, -- $0,439$ против $0,399$ у лучшего лидарного и $0,337$ у визуального -- и сохраняет превосходство во всём диапазоне рассмотренных дальностей. Среднее время обработки кадра $66,7$~мс, что меньше периода регистрации данных $100$~мс.}
\newcommand{\defposB}{Алгоритм проецирования семантической сегментации проезжей части с изображения на оценённую по лидарным данным поверхность движения в сочетании с базовым лидарным алгоритмом повышает качество карты проходимости пространства по сравнению с картой, строящейся исключительно по лидарным данным: среднеквадратичная ошибка стоимости поиска пути (PathFinding Cost Mean Squared Error, PFC-MSE) снижается с $1019,3$ до $924,4$ в пяти рассмотренных сценариях, а среднеквадратичная ошибка относительно эталонной карты -- с $0,231$ до $0,218$. Проецирование той же сегментации на фиксированную плоскость увеличивает PFC-MSE до $1059,4$.}
\newcommand{\defposC}{Структура данных, представляющая локальную карту проходимости пространства набором массивов атомарных пар <<логарифм шансов~-- время>>, делает операции добавления фрагмента, сдвига и получения карты взаимно неблокирующими и уменьшает коэффициент линейной аппроксимации времени итерации цикла публикации карты восемью источниками в $13,5$ раза по сравнению с исходной реализацией.}
\newcommand{\defposD}{Архитектура программного комплекса построения карты проходимости пространства, в которой источник данных представлен поставщиком фрагмента карты, оставляет последовательным только сведение фрагментов в общую карту и сводит подключение датчика нового типа к реализации единственного преобразования его данных во фрагмент. При восьми источниках обновление карты занимает до $22,8$~мс, что меньше периода регистрации данных $100$~мс.}
